\begin{center}
	\large{\textbf{MATRIKS KOEFISIEN BAGI UKURAN KETERBELITAN KEADAAN DUA DAN TIGA QUBIT}}
\end{center}
%\begin{flushleft}
\begin{minipage}{4cm}
	\begin{align*}
		&\text{Nama} \qquad &&: \text{Tamara Pingki} \\
		&\text{NRP} \qquad &&: \text{6001241001} \\
		&\text{Departemen} \qquad &&: \text{Fisika} \\
		&\text{Pembimbing} \qquad &&: \text{Prof. Agus Purwanto, D.Sc.} \\
	\end{align*}
\end{minipage} \\
\begin{center}
	\large\textbf{ABSTRAK} \\  
\end{center}
Keterbelitan kuantum (\textit{quantum entanglement}) merupakan salah satu sumber daya fundamental dalam pemrosesan informasi kuantum. Meskipun klasifikasi keterbelitan pada sistem bipartit telah dipahami dengan baik, identifikasi kelas keterbelitan pada sistem multipartit masih menjadi permasalahan yang kompleks akibat meningkatnya jumlah konfigurasi keterbelitan seiring bertambahnya jumlah qubit. Penelitian ini bertujuan untuk merumuskan suatu kerangka teoritis berbasis matriks koefisien guna mengidentifikasi dan mengklasifikasikan keterbelitan pada sistem dua dan tiga qubit. Metode penelitian dilakukan secara analitik melalui formulasi matriks koefisien yang dibangun dari amplitudo keadaan kuantum. Untuk sistem dua qubit, diperoleh formulasi matriks koefisien berupa $\mathrm{Tr},\rho \tilde{\rho} = 4 \det \rho_A$, sedangkan untuk sistem tiga qubit diperoleh hubungan $\mathrm{Tr}\,\rho_{AB} \tilde{\rho}_{AB}= 4 |\det C_{AB0}|^2 + 4 |\det C_{AB1}|^2 + 2 (\det C_{AB0}^1 + \det C_{AB1}^0)(\det C^{*1}_{AB0} + \det C^{*0}_{AB1})$. Formulasi tersebut kemudian diterapkan pada berbagai keadaan kuantum dua dan tiga qubit untuk menentukan distribusi keterbelitannya. Hasil penelitian ini menunjukkan bahwa pendekatan berbasis matriks koefisien dapat digunakan sebagai metode alternatif yang lebih sederhana dan efisien untuk mengkarakterisasi kelas keterbelitan pada sistem dua dan tiga qubit.

\noindent
\begin{minipage}{1cm}
	\begin{align*}
		\text{\textbf{Kata kunci}}:\hspace{5mm}&\text{Matriks Koefisien, Keterbelitan, Dua Qubit, Tiga Qubit}% \\ &\text{} \\&\text{}
	\end{align*}
\end{minipage}

\newpage
\thispagestyle{empty}
\begin{center}
	\topskip0pt
	\vspace*{\fill}
	\textit{\textbf{"Halaman ini sengaja dikosongkan"}}
	\vspace*{\fill}
\end{center}
\pagebreak